\section{Application Layer}

The application layer provides network service to applications
or users. There are many different application layer protocols,
including but not limited to:
\begin{itemize}
    \item HTTP
    \item TLS
    \item SSL
    \item FTP
    \item SMTP
    \item POP
    \item SSH
    \item NFS
    \item STP
    \item OSPF
    \item RIP
    \item BGP
    \item DNS
\end{itemize}

The World Wide Web is a distributed database of pages linked through
Hyper Text Transfer Protocol. Clients request content served by servers,
while proxies cache, replicate server content for scalability and
availability, better response time, or other reasons. Pages and other
content is identified with a \emph{uniform resource locator} (URL)
and obtained via HTTP.

The syntax of a URL is \texttt{protocol://hostname[:port]/path/resource}.
\begin{itemize}
    \item Protocol: http, ftp, https, etc.
    \item Hostname: DNS name, IP address
    \item Port: default protocol's standard port, e.g. 80 for HTTP, 443 for HTTPS
    \item Path: hierarchical directory path like \texttt{/path/to/directory}
    \item Resource: desired resource, can also execute program
\end{itemize}
For instance, \texttt{htttps://engineering.purdue.edu/~vshriva/index.html}.
